% GENERATED FILE. Edit canonical lessons, then run npm run generate:books.
% canonical-source-hash: fnv1a64:49a673aeb4a74e74
% canonical-lessons: HI-C525-phatakna, HI-C525-ankur, HI-C525-khad, HI-C525-ropna, HI-C525-ukhaarna

\chapter{Seeds and Soil}
\label{ch:hi-seeds-and-soil}
\hlchaptermodality{525}

\begin{chapteropening}
\textbf{By the end of this chapter:} \emph{I can talk about sowing, planting and growing things.}

Complete the last lesson of chapter 525: I can talk about sowing, planting and growing things.
\end{chapteropening}

\section[\texorpdfstring{\hi{फटकना} (phaṭaknā)}{फटकना (phaṭaknā)}]{\hi{फटकना} (phaṭaknā) — to winnow (grain)}

\label{lesson:HI-C525-phatakna}



\begin{quote}
Before the new one: say the Hindi for to set (curd), then the Hindi for to rinse out.
\end{quote}

\subsection*{You'll want to know: \hi{फटकना}}
\textbf{\hi{फटकना}} — \emph{phaṭaknā} — \textquotedblleft{}to winnow (grain)\textquotedblright{}.

\begin{quote}
\hi{किसान} \hi{आँगन} \hi{में} \hi{अनाज} \hi{फटक} \hi{रहा} \hi{है।}
\end{quote}

\begin{quote}
\emph{Kisān āṁgan meṁ anāj phaṭak rahā hai.}
\end{quote}

\begin{quote}
The farmer is winnowing grain in the courtyard.
\end{quote}

Past, with \hi{ने} (he or she …): \hi{उसने} \hi{चावल} \hi{फटके}, \emph{usne cāval phaṭke}.

\begin{cousinweb}
\textbf{\hi{फटकना}} (\emph{phaṭaknā}) is a tadbhava, inherited from \emph{*phaṭ}, \textquotedblleft{}sudden movement\textquotedblright{}, a form CDIAL reconstructs from its descendants (CDIAL 9038). CDIAL gives the Hindi verb as \textquotedblleft{}to winnow\textquotedblright{}.
\end{cousinweb}

\begin{grammarlens}[title={What you've built}]
One more word for describing what you lived and felt.
\end{grammarlens}

\subsection*{Your turn}
\begin{itemize}
\raggedright
  \item \emph{Say it:} \emph{phaṭaknā}
  \item \emph{Say it:} \emph{phaṭaknā}, once more
  \item \emph{Say it:} \emph{Kisān āṁgan meṁ anāj phaṭak rahā hai.} — The farmer is winnowing grain in the courtyard.
  \item \emph{Recall:} say the Hindi for to set (curd), then the Hindi for to rinse out, then say \emph{phaṭaknā} again
\end{itemize}

\begin{tcolorbox}[breakable,colback=teal!4,colframe=teal!35!black,title={Before you move on}]
What does \hi{फटकना} mean? (\textquotedblleft{}To winnow (grain)\textquotedblright{}.) Say it once more.
\end{tcolorbox}
\section[\texorpdfstring{\hi{अंकुर} (aṅkur)}{अंकुर (aṅkur)}]{\hi{अंकुर} (aṅkur) — a sprout; a shoot}

\label{lesson:HI-C525-ankur}



\begin{quote}
Before the new one: say the Hindi for to rinse out, then the Hindi for to winnow (grain).
\end{quote}

\subsection*{You'll want to know: \hi{अंकुर}}
\textbf{\hi{अंकुर}} — \emph{aṅkur} — \textquotedblleft{}a sprout; a shoot\textquotedblright{}.

\begin{quote}
\hi{बीज} \hi{से} \hi{एक} \hi{छोटा} \hi{अंकुर} \hi{निकला।}
\end{quote}

\begin{quote}
\emph{Bīj se ek choṭā aṅkur niklā.}
\end{quote}

\begin{quote}
A small shoot came up from the seed.
\end{quote}

Masculine.

\begin{cousinweb}
\textbf{\hi{अंकुर}} (\emph{aṅkur}) is a tatsama, taken whole from Sanskrit \emph{aṅkura}, \textquotedblleft{}a sprout, a shoot\textquotedblright{} (Monier-Williams).
\end{cousinweb}

\begin{grammarlens}[title={What you've built}]
One more word for describing what you lived and felt.
\end{grammarlens}

\subsection*{Your turn}
\begin{itemize}
\raggedright
  \item \emph{Say it:} \emph{aṅkur}
  \item \emph{Say it:} \emph{aṅkur}, once more
  \item \emph{Say it:} \emph{Bīj se ek choṭā aṅkur niklā.} — A small shoot came up from the seed.
  \item \emph{Recall:} say the Hindi for to rinse out, then the Hindi for to winnow (grain), then say \emph{aṅkur} again
\end{itemize}

\begin{tcolorbox}[breakable,colback=teal!4,colframe=teal!35!black,title={Before you move on}]
What does \hi{अंकुर} mean? (\textquotedblleft{}A sprout; a shoot\textquotedblright{}.) Say it once more.
\end{tcolorbox}
\section[\texorpdfstring{\hi{खाद} (khād)}{खाद (khād)}]{\hi{खाद} (khād) — manure; fertiliser}

\label{lesson:HI-C525-khad}



\begin{quote}
Before the new one: say the Hindi for to winnow (grain), then the Hindi for a sprout.
\end{quote}

\subsection*{You'll want to know: \hi{खाद}}
\textbf{\hi{खाद}} — \emph{khād} — \textquotedblleft{}manure; fertiliser\textquotedblright{}.

\begin{quote}
\hi{खेत} \hi{में} \hi{हर} \hi{साल} \hi{खाद} \hi{डालिए।}
\end{quote}

\begin{quote}
\emph{Khet meṁ har sāl khād ḍāliye.}
\end{quote}

\begin{quote}
Put manure on the field every year.
\end{quote}

Feminine.

\begin{cousinweb}
\textbf{\hi{खाद}} (\emph{khād}) is a tadbhava, inherited from Sanskrit \emph{khāta}, \textquotedblleft{}dug up\textquotedblright{} (CDIAL 3862). CDIAL gives the Hindi word as \textquotedblleft{}manure\textquotedblright{}.
\end{cousinweb}

\begin{grammarlens}[title={What you've built}]
One more word for describing what you lived and felt.
\end{grammarlens}

\subsection*{Your turn}
\begin{itemize}
\raggedright
  \item \emph{Say it:} \emph{khād}
  \item \emph{Say it:} \emph{khād}, once more
  \item \emph{Say it:} \emph{Khet meṁ har sāl khād ḍāliye.} — Put manure on the field every year.
  \item \emph{Recall:} say the Hindi for to winnow (grain), then the Hindi for a sprout, then say \emph{khād} again
\end{itemize}

\begin{tcolorbox}[breakable,colback=teal!4,colframe=teal!35!black,title={Before you move on}]
What does \hi{खाद} mean? (\textquotedblleft{}Manure; fertiliser\textquotedblright{}.) Say it once more.
\end{tcolorbox}
\section[\texorpdfstring{\hi{रोपना} (ropnā)}{रोपना (ropnā)}]{\hi{रोपना} (ropnā) — to plant (a sapling)}

\label{lesson:HI-C525-ropna}



\begin{quote}
Before the new one: say the Hindi for a sprout, then the Hindi for manure.
\end{quote}

\subsection*{You'll want to know: \hi{रोपना}}
\textbf{\hi{रोपना}} — \emph{ropnā} — \textquotedblleft{}to plant (a sapling)\textquotedblright{}.

\begin{quote}
\hi{हमने} \hi{बगीचे} \hi{में} \hi{आम} \hi{का} \hi{एक} \hi{पेड़} \hi{रोपा।}
\end{quote}

\begin{quote}
\emph{Hamne bagīce meṁ ām kā ek peṛ ropā.}
\end{quote}

\begin{quote}
We planted a mango tree in the garden.
\end{quote}

Past, with \hi{ने} (he or she …): \hi{उसने} \hi{फूल} \hi{रोपे}, \emph{usne phūl rope}.

\begin{cousinweb}
\textbf{\hi{रोपना}} (\emph{ropnā}) is a tadbhava, inherited from Sanskrit \emph{rōpayati}, \textquotedblleft{}plants\textquotedblright{}, which CDIAL files with \emph{*rupati}, \textquotedblleft{}pierces, plants\textquotedblright{} (CDIAL 10783). CDIAL gives the Hindi verb as \textquotedblleft{}to place on, plant, sow\textquotedblright{}.
\end{cousinweb}

\begin{grammarlens}[title={What you've built}]
One more word for describing what you lived and felt.
\end{grammarlens}

\subsection*{Your turn}
\begin{itemize}
\raggedright
  \item \emph{Say it:} \emph{ropnā}
  \item \emph{Say it:} \emph{ropnā}, once more
  \item \emph{Say it:} \emph{Hamne bagīce meṁ ām kā ek peṛ ropā.} — We planted a mango tree in the garden.
  \item \emph{Recall:} say the Hindi for a sprout, then the Hindi for manure, then say \emph{ropnā} again
\end{itemize}

\begin{tcolorbox}[breakable,colback=teal!4,colframe=teal!35!black,title={Before you move on}]
What does \hi{रोपना} mean? (\textquotedblleft{}To plant (a sapling)\textquotedblright{}.) Say it once more.
\end{tcolorbox}
\section[\texorpdfstring{\hi{उखाड़ना} (ukhāṛnā)}{उखाड़ना (ukhāṛnā)}]{\hi{उखाड़ना} (ukhāṛnā) — to pull up; to uproot}

\label{lesson:HI-C525-ukhaarna}



\begin{quote}
Before the new one: say the Hindi for manure, then the Hindi for to plant (a sapling).
\end{quote}

\subsection*{You'll want to know: \hi{उखाड़ना}}
\textbf{\hi{उखाड़ना}} — \emph{ukhāṛnā} — \textquotedblleft{}to pull up; to uproot\textquotedblright{}.

\begin{quote}
\hi{माली} \hi{ने} \hi{बगीचे} \hi{से} \hi{पुराना} \hi{पेड़} \hi{उखाड़ा।}
\end{quote}

\begin{quote}
\emph{Mālī ne bagīce se purānā peṛ ukhāṛā.}
\end{quote}

\begin{quote}
The gardener pulled up the old tree from the garden.
\end{quote}

Past, with \hi{ने} (he or she …): \hi{उसने} \hi{घास} \hi{उखाड़ी}, \emph{usne ghās ukhāṛī}.

\begin{cousinweb}
\textbf{\hi{उखाड़ना}} (\emph{ukhāṛnā}) is a tadbhava, inherited from \emph{*utskarati}, \textquotedblleft{}digs out\textquotedblright{}, a form CDIAL reconstructs from its descendants (CDIAL 1895). CDIAL gives the Hindi verb as \textquotedblleft{}to uproot\textquotedblright{}.
\end{cousinweb}

\begin{grammarlens}[title={What you've built}]
One more word for describing what you lived and felt.
\end{grammarlens}

\subsection*{Your turn}
\begin{itemize}
\raggedright
  \item \emph{Say it:} \emph{ukhāṛnā}
  \item \emph{Say it:} \emph{ukhāṛnā}, once more
  \item \emph{Say it:} \emph{Mālī ne bagīce se purānā peṛ ukhāṛā.} — The gardener pulled up the old tree from the garden.
  \item \emph{Recall:} say the Hindi for manure, then the Hindi for to plant (a sapling), then say \emph{ukhāṛnā} again
\end{itemize}

\begin{tcolorbox}[breakable,colback=teal!4,colframe=teal!35!black,title={Before you move on}]
What does \hi{उखाड़ना} mean? (\textquotedblleft{}To pull up; to uproot\textquotedblright{}.) Say it once more.
\end{tcolorbox}
